\documentclass[11pt]{article}
\usepackage[margin=1in]{geometry}
\usepackage{hyperref}
\usepackage{longtable}
\usepackage{array}
\usepackage{enumitem}

\title{Intelligent Attendance System Using Face Recognition\\Project Report}
\author{Project Review}
\date{April 5, 2026}

\begin{document}
\maketitle

\section*{Executive Summary}
This project is a local, academic attendance system that uses face detection and face recognition to mark attendance automatically. It combines a Python FastAPI backend for capture, training, recognition, and data storage with a Next.js frontend for administration and reporting. All data is stored locally in SQLite, and face images are stored on disk.

\section*{Goals and Scope}
\begin{itemize}[leftmargin=*]
  \item Automate daily attendance using a webcam and face recognition.
  \item Provide simple admin screens for student registration, model training, and attendance reporting.
  \item Keep the system local and lightweight for classroom or lab demos.
\end{itemize}

\section*{Technology Stack}
\begin{itemize}[leftmargin=*]
  \item Backend: Python 3.9+, FastAPI, SQLAlchemy, OpenCV, NumPy, Pandas, Pillow.
  \item Face recognition: Haar Cascade for detection, LBPH for recognition.
  \item Frontend: Next.js 16, React 19, Tailwind CSS 4.
  \item Database: SQLite (local file).
\end{itemize}

\section*{System Architecture}
\begin{itemize}[leftmargin=*]
  \item The frontend calls the backend over HTTP at \texttt{http://localhost:8000}.
  \item The backend manages webcam capture, model training, recognition, and persistence.
  \item Attendance is written to SQLite; face images and the trained model are stored on disk.
\end{itemize}

\section*{Directory Map (Observed)}
\begin{verbatim}
.
|-- project_plan.md
|-- README.md
|-- attendance-system/
|   |-- backend/
|   |   |-- attendance_service.py
|   |   |-- database.py
|   |   |-- dataset/ (User.<id>.<count>.jpg)
|   |   |-- face_detection.py
|   |   |-- face_recognition.py
|   |   |-- haarcascade_frontalface_default.xml
|   |   |-- main.py
|   |   |-- models.py
|   |   |-- register_student.py
|   |   |-- static/ (empty)
|   |   |-- train_model.py
|   |   |-- trainer/trainer.yml
|   |-- database/attendance.db
|-- frontend/
|   |-- package.json
|   |-- next.config.ts
|   |-- src/app/
|   |   |-- layout.tsx
|   |   |-- page.tsx
|   |   |-- attendance/page.tsx
|   |   |-- students/page.tsx
|   |   |-- reports/page.tsx
|   |   |-- globals.css
|   |-- public/ (svg assets)
\end{verbatim}

\section*{Backend Design}
\subsection*{Core Modules}
\begin{itemize}[leftmargin=*]
  \item \texttt{main.py}: FastAPI entry point; defines all API routes and camera streaming.
  \item \texttt{database.py}: SQLAlchemy engine, session, and DB path configuration.
  \item \texttt{models.py}: Student and Attendance models; creates tables on import.
  \item \texttt{register\_student.py}: Captures 30 face images for a student ID.
  \item \texttt{train\_model.py}: Trains the LBPH recognizer from captured images.
  \item \texttt{face\_detection.py}: Haar Cascade detection; returns face boxes and grayscale frame.
  \item \texttt{face\_recognition.py}: Loads LBPH model and predicts face ID.
  \item \texttt{attendance\_service.py}: Attendance marking and CSV export.
\end{itemize}

\subsection*{API Endpoints}
\begin{longtable}{|p{0.24\textwidth}|p{0.12\textwidth}|p{0.57\textwidth}|}
\hline
\textbf{Route} & \textbf{Method} & \textbf{Purpose} \\\hline
\texttt{/} & GET & Health check; returns a startup message. \\\hline
\texttt{/register-student} & POST & Creates a student, then opens webcam to capture 30 images. \\\hline
\texttt{/train-model} & POST & Trains LBPH and saves \texttt{trainer.yml}. \\\hline
\texttt{/attendance} & GET & Returns attendance records joined with student data. \\\hline
\texttt{/export-csv} & GET & Streams CSV export of attendance records. \\\hline
\texttt{/students} & GET & Returns all students. \\\hline
\texttt{/delete-student/\{id\}} & DELETE & Deletes student, their attendance, and images. \\\hline
\texttt{/toggle-camera} & POST & Starts or stops the camera stream (boolean query). \\\hline
\texttt{/video-feed} & GET & MJPEG stream when camera is active. \\\hline
\texttt{/start-attendance} & GET & Opens a local CV2 window for recognition. \\\hline
\end{longtable}

\section*{Face Recognition Pipeline}
\begin{enumerate}[leftmargin=*]
  \item Registration: \texttt{capture\_student\_images} stores 30 grayscale face crops in \texttt{dataset/}.
  \item Training: \texttt{train\_lbph\_model} reads \texttt{User.<id>.<count>.jpg} files and saves \texttt{trainer.yml}.
  \item Recognition: \texttt{detect\_faces} finds faces; \texttt{recognize\_face} predicts ID; \texttt{mark\_attendance} records once per day.
\end{enumerate}

\section*{Data Model}
\begin{longtable}{|p{0.22\textwidth}|p{0.70\textwidth}|}
\hline
\textbf{Table} & \textbf{Fields} \\\hline
\texttt{students} & \texttt{id} (PK), \texttt{name}, \texttt{roll\_number} (unique) \\\hline
\texttt{attendance} & \texttt{id} (PK), \texttt{student\_id} (FK), \texttt{date}, \texttt{time}, \texttt{status} \\\hline
\end{longtable}

\section*{Frontend Design}
\begin{itemize}[leftmargin=*]
  \item Global layout provides navigation: Dashboard, Students, Start Attendance, Reports.
  \item Dashboard shows counts and triggers model training.
  \item Students page registers new students and deletes existing students.
  \item Attendance page starts or stops live recognition and shows MJPEG feed.
  \item Reports page lists attendance and provides CSV export.
\end{itemize}

\section*{Data Flow (End-to-End)}
\begin{enumerate}[leftmargin=*]
  \item Admin registers a student in the UI; backend creates record and captures images.
  \item Admin trains the model; LBPH model is saved to \texttt{trainer/trainer.yml}.
  \item Admin starts attendance; backend recognizes faces and writes attendance rows.
  \item Reports page fetches data and exports CSV on demand.
\end{enumerate}

\section*{Runtime and Operations}
\begin{itemize}[leftmargin=*]
  \item Backend: run \texttt{python main.py} in \texttt{attendance-system/backend}.
  \item Frontend: run \texttt{npm install} then \texttt{npm run dev} in \texttt{frontend}.
  \item Backend listens on \texttt{localhost:8000}, frontend on \texttt{localhost:3000}.
\end{itemize}

\section*{Current Artifacts Observed}
\begin{itemize}[leftmargin=*]
  \item SQLite file: \texttt{attendance-system/database/attendance.db}.
  \item Training output: \texttt{attendance-system/backend/trainer/trainer.yml}.
  \item Sample dataset: \texttt{attendance-system/backend/dataset/User.1.*.jpg}.
\end{itemize}

\section*{Risks and Limitations}
\begin{itemize}[leftmargin=*]
  \item LBPH accuracy depends on lighting, pose, and camera quality.
  \item No authentication or authorization on API endpoints.
  \item Model and images are stored locally; no encryption or backup.
  \item Web streaming is MJPEG and may not scale for many clients.
\end{itemize}

\section*{Suggested Enhancements}
\begin{itemize}[leftmargin=*]
  \item Add admin login and role-based access.
  \item Move host and port settings to environment variables.
  \item Add unit tests for API endpoints and attendance logic.
  \item Improve recognition with more robust models and thresholds.
  \item Add pagination and filters on the reports page.
\end{itemize}

\end{document}
